\title{Instructions for how to use this \texttt{polarimetryLab} program}
\author{Yakir Gagnon}
\date{\today}
\documentclass[12pt]{article}
\usepackage{graphicx}
\begin{document}
\maketitle

\section{Introduction}
This program computes polarimetry images.

\paragraph{Outline}
The program takes the six filtered images (i.e.\ anti-diagonal, vertical, diagonal, horizontal, lefthand circular, and righthand circular) and produces five polarization images that describe the RGB, intensity, degree of linear polarization, angle of polarization, and degree of circular polarization. 

\section{Running the program}
\begin{enumerate}
	\item Photograph your subject through the following filters in the following order: anti-diagonal, vertical, diagonal, horizontal, and if you want circular polarization as well: lefthand circular, and righthand circular (see Section \ref{shooting} for detailed instructions for how to shoot these images). 
	\item Put (copy/move) these images in the folder \texttt{raw}, making sure you have only 4 or 6 images (depending if you want circular polarization or not) in the \texttt{raw} folder.
	\item Double click the icon \includegraphics[height=2ex]{icon}.
	\item Wait for the "Done" message.
	\item The polarimetry images are now in the folder \texttt{polarimetry}.
	\item Take (copy/move) the resulting images from the \texttt{polarimetry} folder because any subsequent runs will erase anything in the \texttt{polarimetry} folder.
\end{enumerate}

\section{Photographing}\label{shooting}
For the best results use the D300 Nikon with the filter rings attached to the lens. 

\end{document}
